\section{Proofs}
\label{app:proofs}

Throughout, $\PP_x,\E_x$ condition on $X{=}x$; the lag $T$ is independent of reading errors given $X$; and the policy's randomization $U$ has a law that does not depend on $X$.
For Theorem~\ref{thm:main} the family's reads are conditionally i.i.d.\ given $X$ and the visibility state, with $Z_i=\log\frac{p_1(Y_i)}{p_0(Y_i)}$ and $L_n=\sum_{i\le n}Z_i$.
Because $p_0\ne p_1$ have full support, $\E_1Z=D_1>0$, $\E_0Z=-D_0<0$ and $\PP_x(Z\ne0)>0$.

\subsection{Theorem~\ref{thm:main}}
The theorem in full reads as follows. For a family of conditionally i.i.d.\ views with full support, let $u^\star=\max_y\log\frac{p_1(y)}{p_0(y)}$, $v^\star=\max_y\log\frac{p_0(y)}{p_1(y)}$
and $d(x\|y)=x\log\frac{x}{y}+(1-x)\log\frac{1-x}{1-y}$. Consider an effect with prior $\pi\in(0,1)$, lag tail $G$, and views with laws $p_0,p_1$.
\begin{enumerate}[leftmargin=*,label=(\alph*),topsep=1pt,itemsep=1pt]
\item \emph{Achievability.} If $G(w)<\alpha$, waiting until $r_e+w$ and then running the SPRT on $L_n=\sum_{i\le n}\log\frac{p_1(Y_i)}{p_0(Y_i)}$, which declares when
$L_n\ge a=\log\frac1\beta$ and retries when $L_n\le-b=-\log\frac{1}{\alpha-G(w)}$, meets $(\alpha,\beta)$. Given visibility,
\begin{equation}
\E_1[N]\le\frac{a+u^\star}{D_1},\qquad \E_0[N]\le\frac{b+v^\star}{D_0}.
\label{eq:ub}
\end{equation}
\item \emph{Lag converse.} Every policy that decides by $r_e+w$ from end-to-end reads alone has $\PP(\text{retry}\mid X{=}1)\ge G(w)\,\PP(\text{retry}\mid X{=}0)$.
A verification-only one with $\PP(\text{declare}\mid X{=}0)\le\beta$ thus has $\PP(\text{retry}\mid X{=}1)\ge(1-\beta)\,G(w)$.
\item \emph{Information converse.} Every verification-only policy meeting $(\alpha,\beta)$ with $\alpha+\beta<1$ and $\E_x[N]<\infty$ satisfies the right inequality
below for any lag, and the left one if all reads follow visibility:
\begin{equation}
\E_1[N]\ge\frac{d(1-\alpha\,\|\,\beta)}{D_1},\qquad \E_0[N]\ge\frac{d(1-\beta\,\|\,\alpha)}{D_0}.
\label{eq:lb}
\end{equation}
\item \emph{Weak evidence.} Every policy that does not ask has $\PP(\dup)+\PP(\loss)\ge\E[\min(\pi_{\mathcal F},1-\pi_{\mathcal F})]$, where $\pi_{\mathcal F}$ is the posterior
of $X{=}1$ when it decides, and the Bayes rule attains this. Under any lag, the right side is at least $\min(\pi,1-\pi)\big(1-\sqrt{D_0\E_0[N]/2}\big)$.
\end{enumerate}
The gap between~\eqref{eq:ub} and~\eqref{eq:lb} under $X{=}1$ tends to $u^\star/D_1$ as $\alpha,\beta\to0$, provided $\alpha\log(1/\beta)\to0$. A fallible user makes asking cost
$c_{\mathrm{ask}}+\epsilon_u[\pi_{\mathcal F}c_{\dup}+(1-\pi_{\mathcal F})c_{\loss}]$.

\para{(a)} Let $N=\inf\{n:L_n\notin(-b,a)\}$ over reads taken after $r_e+w$. Since $\PP_x(Z\neq0)>0$, Stein (1946) gives $\PP_x(N>n)\le cr^n$ for constants $c<\infty$, $r<1$, so
$\E_x[N]<\infty$~\cite{tartakovsky2014sequential}. Under $X{=}0$ reads are i.i.d.\ $p_0$, and the likelihood-ratio identity for stopped sequences gives
$\PP_0(L_N\ge a)=\E_1[e^{-L_N};L_N\ge a]\le e^{-a}=\beta$. Under $X{=}1$,
$\PP_1(\text{retry})\le\PP_1(T>w)+\PP_1(\text{retry}\mid T\le w)$. Given $T\le w$ every read is i.i.d.\ $p_1$, so $\PP_1(L_N\le-b\mid T\le w)=\E_0[e^{L_N};L_N\le-b]\le e^{-b}=\alpha-G(w)$.
Both inequalities are exact; overshoot only lowers the realized error. For~\eqref{eq:ub}, Wald's identity gives $\E_1[L_N]=D_1\E_1[N]$ given visibility; on $\{L_N\ge a\}$,
$L_{N-1}<a$ so $L_N<a+u^\star$, and on $\{L_N\le-b\}$, $L_N<a$. Hence $D_1\E_1[N]\le a+u^\star$; symmetrically $D_0\E_0[N]\le b+v^\star$.

\para{(b)} Let the decision be $\phi(U,Y_1,\dots,Y_N)$ with all reads end-to-end and taken by $r_e+w$. Given $X{=}1$ and $T>w$, A1 says every such read has the law it has under $X{=}0$,
and $U$ is independent of $(X,T)$. So $(U,N,Y_{1:N})$ has the same law as under $X{=}0$, $\PP_1(\text{retry}\mid T>w)=\PP_0(\text{retry})$, and
$\PP_1(\text{retry})\ge G(w)\,\PP_0(\text{retry})$. A verification-only policy has $\PP_0(\text{retry})=1-\PP_0(\text{declare})\ge1-\beta$. The argument fails for a policy that may ask
($\PP_0(\text{retry})$ can be 0) or reads a hop-1 view (whose law does not depend on $T$).

\para{(c)} Let $Q_x$ be the law of $(U,N,Y_1,\dots,Y_N)$ under $\PP_x$. Without lag, the log-likelihood ratio of $Q_1$ to $Q_0$ is $L_N$ (the $U$ terms cancel), so Wald's identity gives
$\KL(Q_1\|Q_0)=D_1\E_1[N]$. With a lag, $Q_1=\int Q_1^{t}\,\mu(dt)$ is a mixture over $T$; under $Q_1^t$ reads before visibility have law $p_0$ and contribute nothing to
$\KL(Q_0\|Q_1^t)$, so $\KL(Q_0\|Q_1^t)\le D_0\E_0[N]$, and by convexity of KL in its second argument $\KL(Q_0\|Q_1)\le D_0\E_0[N]$. The decision is a function of
$(U,N,Y_{1:N})$, so by data processing~\cite{cover2006elements}
\[
\KL(Q_0\|Q_1)\ge d\big(\PP_0(\text{retry})\,\|\,\PP_1(\text{retry})\big)=d(1-\beta'\|\alpha'),
\]
with $\alpha'\le\alpha$, $\beta'\le\beta$. For $p>q$, $d(p\|q)$ increases in $p$ ($\partial_pd=\log\frac{p(1-q)}{q(1-p)}>0$) and decreases in $q$ ($\partial_qd=\frac{q-p}{q(1-q)}<0$); since
$1-\beta'\ge1-\beta>\alpha\ge\alpha'$, $d(1-\beta'\|\alpha')\ge d(1-\beta\|\alpha)$. The bound on $\E_1[N]$ is symmetric without lag~\cite{kaufmann2016complexity}.

\para{(d)} Let $r_{\mathcal F}=\PP(\text{retry}\mid\mathcal F)$. Then $\PP(\dup)+\PP(\loss)=\E[\pi_{\mathcal F}r_{\mathcal F}+(1-\pi_{\mathcal F})(1-r_{\mathcal F})]\ge\E[\min(\pi_{\mathcal F},1-\pi_{\mathcal F})]$,
with equality for the Bayes rule. For the second claim, $\pi\,Q_1(\text{retry})+(1-\pi)Q_0(\text{declare})\ge\min(\pi,1-\pi)[1-(Q_0(\text{retry})-Q_1(\text{retry}))]\ge\min(\pi,1-\pi)(1-\mathrm{TV}(Q_0,Q_1))$,
and Pinsker gives $\mathrm{TV}\le\sqrt{\KL(Q_0\|Q_1)/2}\le\sqrt{D_0\E_0[N]/2}$ by~(c), under any lag.

\para{Corollary~\ref{cor:notfree}} (absolute continuity). If $\PP_0(\text{retry})>0$, some finite record $(u,y_{1:n})$, with reads at times $t_1\le\dots\le t_n$, leads to retry and has
positive $\PP_0$-probability. If $0<\epsilon_1<1$, $p_1$ has full support and the record has positive $\PP_1$-probability, so $\PP_1(\text{retry})>0$. If $G(t)>0$ for all $t$, then under
$X{=}1$ and $T>t_n$ every read has its $X{=}0$ law (A1), so $\PP_1(\text{retry})\ge G(t_n)\,\PP_0(\text{record})>0$. Neither argument needs a bounded stopping time.
For the equivalence: if $D_0>0$, pick $w$ with $G(w)<\alpha$ ($G\to0$ since $X=\mathbf 1\{T<\infty\}$) and apply (a), which needs $\beta>0$ for a finite threshold $a$. If $D_0=0$ then $p_0=p_1$,
every read has the same law under both outcomes and at all times, so $\PP_1(\text{retry})=\PP_0(\text{retry})\ge1-\beta>\alpha$.

\subsection{Product certification}
\begin{proposition}[Family-level certification]
\label{prop:family}
Let the planner consult families in a fixed order, multiply the bounds of the verdicts seen so far, retry at the first point where the product of
$q^{\mathrm{abs}}$ terms is at most $\tau$, and ask on a contradiction. Under A2,
\[
\PP(\text{retry}\mid X{=}1,\text{visible})\;\le\;B(\tau)=\sum_{A\in\mathcal R(\tau)}\;\prod_{(f,r)\in A}q^{\mathrm{abs}}_f(r),
\]
where $\mathcal R(\tau)$ is the set of ordered assignments $A$ (each family skipped or at a level $r$) with $\prod_A q\le\tau$ and $\prod_{A\setminus\mathrm{last}(A)}q>\tau$.
The same holds for declaring with $q^{\mathrm{com}}$. No assumption is made about dependence among the views of a family.
If a hop-1 entry ends failed or stuck, the planner hands over: it resets the product to $1$ exactly once per effect and continues on the end-to-end families.
$\mathcal R(\tau)$ includes the assignments that start after the hand-over, which skip hop~1, so the bound covers this path.
\end{proposition}
\begin{proof}
Fix the family order $f_1,\dots,f_m$. A run of the planner that retries stops right after some family $f_j$, having seen abs-verdicts at levels $r$ for a subset of
$f_1,\dots,f_j$ (the others tied, i.e.\ skipped) and no com-verdict (a com-verdict before or after would be a contradiction, which asks). The ordered list $A$ of these
(family, level) pairs satisfies $\prod_Aq\le\tau$, and $\prod_{A\setminus(f_j,r_j)}q>\tau$ because the planner did not stop earlier. So $A\in\mathcal R(\tau)$, and the run lies in
$E_A=\bigcap_{(f,r)\in A}\{\text{at least $r$ views of $f$ read absent}\}$. Under $X{=}1$ with the effect visible, A2 makes the events in $E_A$ independent, so
$\PP(E_A)\le\prod_{(f,r)\in A}q^{\mathrm{abs}}_f(r)$. A union bound over $\mathcal R(\tau)$ gives $B(\tau)$. Nothing is assumed about the joint law of views within a family,
because each $q$ bounds the family's verdict event directly.
For the hand-over: if hop~1 ends failed or stuck, the planner sets the product to $1$ once and does not use the hop-1 verdict again. A run that then retries stops after an ordered
list of end-to-end verdicts that skips hop~1. That list satisfies the two conditions above, so it is in $\mathcal R(\tau)$ and the same union bound covers it. The reset happens at most
once per effect, so no run is counted under two lists.
\end{proof}


\emph{Lag, Eq.~\eqref{eq:lagsplit}.} Fix the class $c$ and write $T=T_c$. Since $t_{\mathrm{vis}}\le r_e+T$, a read at time $t\ge r_e+w$ on the event $T\le w$ is post-visibility. Partition retrying runs by path: $\Pi_1$, hop~1 read ``absent'' (end-to-end reads from $r_e+w_1$ on); $\Pi_0$, every other path (end-to-end reads
from $r_e+w_0$ on, measured from $r_e$ also after a stuck or failed hand-over). On $\{T\le w_1\}\cap\Pi_1$ or $\{T\le w_0\}\cap\Pi_0$ every end-to-end read is post-visibility, and the run
lies in some $E_A$, $A\in\mathcal R(\tau)$, so these parts sum to at most $B(\tau)$ (the paths are disjoint, and assignments that skip hop~1 belong to $\mathcal R(\tau)$).
$\PP(\Pi_1,T>w_1\mid X{=}1)=\PP(\text{hop-1 abs}\mid X{=}1)\,G(w_1)\le q_hG(w_1)$, since hop-1 readings and $T$ are independent given $X$. $\PP(\Pi_0,T>w_0\mid X{=}1)\le G(w_0)$,
since $\Pi_0$ is a single event set; on the stuck path hop-1 status depends on $T$, which this term does not need. The bound $q_hG(w)\le\alpha_{\mathrm{lag}}$ alone would spend the lag
share twice when both paths occur (up to $(2-q_h)\alpha_{\mathrm{lag}}$).

\subsection{Proposition~\ref{prop:joint}}
The regions $R_{\tau_1}\subseteq R_{\tau_2}\subseteq\cdots$ are fixed by the pilot split, so their true probabilities $r_1\le r_2\le\cdots$ are non-decreasing. If every $r_j\le\alpha_{\mathrm{test}}$
there is nothing to show. Otherwise let $j^\star$ be the first index with $r_{j^\star}>\alpha_{\mathrm{test}}$. The procedure keeps some $\tau_j$ with $j\ge j^\star$ only if the CP bound for $R_{\tau_{j^\star}}$
is at most $\alpha_{\mathrm{test}}<r_{j^\star}$, i.e.\ it undercovers, which has probability at most $\delta$. A retry at threshold $\tau$ in any configuration puts the realized tuple (all reads post-visibility) in $R_\tau$, which unites the configurations; this gives
the claim. The CP bound with zero events is $1-\delta^{1/n}$, which is at most $\alpha_{\mathrm{test}}$ only if $n\ge\ln(1/\delta)/\alpha_{\mathrm{test}}$ (approximately).

For source families, the mode keeps A2 only between hop~1, which reads the app's Outbox, and the other three families, which read the server-rendered history.
It is available only when hop~1 is read from a store distinct from the rendered history. We check this per app by testing hop-1 readings against $T$ on fault-injected
trials, and Element fails the check. Merging the history readers into one jointly calibrated family gives $B(\tau)=q_h\,Q(\tau/q_h)+Q(\tau)$. Here $Q(t)$ is a CP bound on the
history family reaching product $t$, with $Q(t)=0$ below the smallest reachable product and $Q(t)=1$ for $t\ge1$ (hop~1 alone then retries, whatever the history shows).
The proof applies Proposition~\ref{prop:family} to the two families \{hop~1\} and \{history readers\}. Both $Q$ terms are tested in the same fixed sequence, so the failure
probability is at most $2\delta$.

\subsection{Persistent-error floor}
\begin{lemma}[Persistent-error floor]
\label{lem:floor}
Let the $V$ views of a family err conditionally i.i.d.\ given a latent per-effect error rate $\theta$ with $\E\theta=\epsilon$ and $\mathrm{Var}\,\theta=\rho\epsilon(1-\epsilon)$,
so that two views' errors have correlation $\rho$. Then
\begin{gather*}
\sup\PP(\text{all $V$ views wrong})=w+(1-w)\big(\epsilon(1-\rho)\big)^V,\\
w=\rho\epsilon/(1-\epsilon+\rho\epsilon),
\end{gather*}
attained when the family is \emph{blind} (all views wrong) with probability $w$ and otherwise errs independently at rate $\epsilon(1-\rho)$.
If the views are only exchangeable, with $m_2=\epsilon^2+\rho\epsilon(1-\epsilon)$, the supremum is $m_2$ for $2\le V\le(1-m_2)/(\epsilon-m_2)$ and strictly smaller for larger $V$.
\end{lemma}
As $V\to\infty$ the supremum tends to $w>0$, so more views of one family cannot certify a side below its persistent mass. A hidden nonce field is such a blind state ($w\ge\eta$),
so certifying ``absent'' at $\alpha<\eta$ needs a second family. Within-view correlation does not matter, because no view is read twice.

\begin{proof}
Let $a=\epsilon(1-\rho)$ and $f(\theta)=\theta^V$. For $V\le2$, $\E f(\theta)$ is fixed by the first two moments, so every admissible law, including the two-point law below, gives the same
value. For $V\ge3$, let $q$ be the quadratic with $q(a)=f(a)$, $q'(a)=f'(a)$ and $q(1)=f(1)$. The Hermite interpolation error is $f(\theta)-q(\theta)=\tfrac{1}{6}f'''(\xi)(\theta-a)^2(\theta-1)$
for some $\xi\in(0,1)$; since $f'''\ge0$ on $[0,1]$, $f\le q$ there. So $\E f(\theta)\le\E q(\theta)$, and $\E q(\theta)$ depends only on $\E\theta$ and $\E\theta^2$. The two-point law
that puts mass $w=\rho\epsilon/(1-\epsilon+\rho\epsilon)$ on $1$ and $1-w$ on $a$ has mean $\epsilon$ and variance $\rho\epsilon(1-\epsilon)$ (direct computation), and $f=q$ on its support,
so it attains $\E q$. Hence $\sup\E f(\theta)=w+(1-w)a^V$, and $\PP(\text{all $V$ wrong})=\E\theta^V$ by conditional independence. As $V\to\infty$ the supremum decreases to $w$.
\emph{Exchangeable views.} Let $\xi_1,\dots,\xi_V$ be exchangeable error indicators with mean $\epsilon$ and pairwise correlation $\rho$, and $m_2=\E\xi_1\xi_2=\epsilon^2+\rho\epsilon(1-\epsilon)$.
Then $\PP(\text{all wrong})\le\PP(\xi_1=\xi_2=1)=m_2$. The law of $S=\sum\xi_i$ with $\PP(S{=}V)=m_2$, $\PP(S{=}1)=V(\epsilon-m_2)$ and $S=0$ otherwise (each realization's positions uniform)
has marginal $\epsilon$ and pair probability $m_2$, and is valid when $V(\epsilon-m_2)\le1-m_2$, where it attains $m_2$. For larger $V$ the mass $V(\epsilon-m_2)$ that
must sit on configurations with exactly one error exceeds $1-m_2$, so any law with pair probability $m_2$ puts pair mass on configurations with $2\le S<V$, and
$\PP(S{=}V)<m_2$. For $V{=}1$ the probability is $\epsilon$.
\end{proof}


\subsection{Lemma~\ref{lem:calib}}
(i)~Let $\tilde u=\log\frac{1-\tilde\epsilon_1}{\tilde\epsilon_0}$, $\tilde v=\log\frac{1-\tilde\epsilon_0}{\tilde\epsilon_1}$, both positive. Under $X{=}0$, drive the true reads and the model's
reads with the same uniforms $U_i$: a true read is 1 iff $U_i<\epsilon_0$, a model read iff $U_i<\tilde\epsilon_0$. Since $\epsilon_0\le\tilde\epsilon_0$, $L^{\mathrm{true}}_n\le L^{\mathrm{model}}_n$
for all $n$. If the true walk exits at the top at step $n$, the model walk stayed above $-b$ before $n$ and is at least $a$ at $n$, so it exits at the top no later. Hence
$\PP^{\mathrm{true}}_0(\text{declare})\le\PP^{\mathrm{model}}_0(\text{declare})\le\beta$ by (a) applied to the model, under which the planner's statistic is the exact log-likelihood ratio.
Under $X{=}1$ a true 0-read implies a model 0-read, and the same argument bounds retries by $\alpha$. The argument is per read, so it also covers independent, non-identical views.
(ii)~A one-sided $(1-\delta)$ Clopper--Pearson bound covers its parameter with probability at least $1-\delta$~\cite{clopper1934use} when the calibration effects are i.i.d.\ from the deployment
population of the effect's fault class; a union bound over the $Kn_{\mathrm{cls}}$ bounds of all classes gives $1-Kn_{\mathrm{cls}}\delta$ (for an unknown class, the per-class bounds that the fallback uses are all in the union, and
a per-class bound valid for every class bounds any mixture of classes), and on the complementary event the guarantee may fail with probability at most 1, which gives the marginal bound. The
declare side is identical with the ``committed'' bounds and $\beta_{\mathrm{test}}$.

\subsection{Drift}
\begin{lemma}[Independent drift]
\label{lem:drift}
If deployed family rates exceed their calibrated bounds by factors $\kappa_f\ge1$, the test term of~\eqref{eq:lagsplit} is at most $\max_{A\in\mathcal R(\tau)}\prod_{f\in A}\kappa_f\cdot B(\tau)$.
\end{lemma}
\begin{proof}
Each term of $B(\tau)$ is $\prod_{f\in A}q_f$. Replacing $q_f$ by $\kappa_fq_f$ multiplies it by $\prod_{f\in A}\kappa_f$, at most the maximum over $\mathcal R(\tau)$.
\end{proof}
The lemma is a margin against drift that is independent across families. Common-mode drift breaks the product itself, which is what joint certification is for.

\subsection{Content-carried nonce}
\begin{lemma}[Content-carried nonce]
\label{lem:nonce}
Suppose the effect has a free-text field that the app stores verbatim and renders in sent/history views, and that the probe reads only those views, only outgoing items of
the verified account, and never drafts, failed entries, notifications or keyboard suggestions, so an uncommitted nonce is never visible to it. The planner writes $k$ symbols
drawn uniformly from an alphabet of size $m$ and reuses them on retry. A view matches if one of its $M$ visible tokens is within edit distance $t$ of the nonce after OCR in which
each symbol errs independently with probability $e_c$. If the field is hidden, truncated or mis-tokenized (split, merged, wrapped) with probability $\eta$, then per view
\begin{equation}
\epsilon_1\le\eta+\PP[\mathrm{Bin}(k,e_c)>t],\qquad \epsilon_0\le M\,|B_t|/m^k,
\label{eq:nonce}
\end{equation}
where $|B_t|$ bounds the number of length-$k$ strings within edit distance $t$ of any token ($|B_0|=1$, $|B_1|\le km$), regardless of look-alike effects on screen.
\end{lemma}
The hidden-field term $\eta$ is a blind state in the sense of Lemma~\ref{lem:floor}, common to all nonce views.

\begin{proof}
\emph{Miss.} If $X{=}1$ and the effect is visible, the nonce token is on screen unless the field is hidden or truncated (probability $\le\eta$). Each symbol suffers at most one OCR
operation with probability $e_c$, so the edit distance between the read token and the nonce is at most a $\mathrm{Bin}(k,e_c)$ count; a union bound gives $\epsilon_1$.
\emph{False match.} If $X{=}0$, the nonce was never committed and the probe never reads a view where an uncommitted copy could appear, so the $M$ visible tokens and their readings are
independent of the uniformly drawn nonce $\nu$. For a fixed reading $s$, $\PP(d(s,\nu)\le t)\le|B_t|/m^k$, and a union bound over $M$ tokens gives $\epsilon_0$. For $t{=}1$, the
length-$k$ strings within edit distance 1 of $s$ are: $s$ and its $k(m-1)$ substitutions if $|s|=k$; at most $km$ insertions if $|s|=k-1$; at most $k+1$ deletions if $|s|=k+1$;
none otherwise. All are at most $km$.
\end{proof}


\subsection{Theorem~\ref{thm:twohop}}
Let hop~1 create $n_1$ entries with distinct ids. By H1 every application at the server is of an outbox entry's id.
($\Leftarrow$)~With $n_1{=}1$ and G2: durability keeps the entry across admissible app kills (which do not cancel scheduled work), and retransmission until acknowledged, with finitely many
faults, ends with an applied, acknowledged transmission, so $C_e\ge1$. By H2 every retransmission carries the same id, and the server applies it at most once, so $C_e\le1$.
($\Rightarrow$, contrapositive)~If $n_1{=}0$, nothing is sent and $C_e=0$ on the fault-free schedule. If $n_1\ge2$, the entries carry distinct ids, dedup does not merge them, and
$C_e\ge2$ on the fault-free schedule. If $n_1{=}1$ and G2 fails: without durability, ``app kill after hop~1, before the first transmission'' gives $C_e=0$; without retransmission,
``first transmission lost before apply'' gives $C_e=0$; without dedup, ``first transmission applied, acknowledgement lost, retransmission applied'' gives $C_e=2$. Each schedule is admissible.
Under G2$'$, dedup still gives $C_e\le1$; retransmission stops before an acknowledgement only in the failed state, so $C_e=0$ implies the entry ended failed.
\emph{Horizon.} If a retransmission may come later than $H_{\mathrm{dedup}}$ after the first application, ``first transmission applied, acknowledgement lost, retransmission after $H_{\mathrm{dedup}}$'' gives $C_e=2$,
so the horizon condition in H3 is necessary; Synapse showed exactly this after we aged a transaction-id row past 24\,h (\S\ref{sec:eval-real}).

Dedup scopes differ by protocol. Matrix clients send with a client-chosen transaction id that makes retransmission idempotent~\cite{matrix_spec}. Synapse keeps
these ids per device in a 30--60\,min cache and a 24\,h table~\cite{synapse_txn,matrix_spec}, and we measured both scopes (\S\ref{sec:eval-real}). CalDAV creates resources
at a client-chosen URL with a conditional \texttt{PUT}~\cite{rfc4791,rfc9110}. SMTP has no end-to-end dedup~\cite{rfc5321}. Element and DAVx5 resolved every force-stop
by relaunch, although force-stop violates H3.

\subsection{Proposition~\ref{prop:saga}}
For a workflow with several ambiguous effects~\cite{garciamolina1987sagas,chang2025sagallm}:
\begin{proposition}[Saga budget allocation]
\label{prop:saga}
Per-step guarantees $(\alpha_i,\beta_i)$ give $\PP(\text{any dup})\le\sum_i\pi_i\alpha_i\le\sum_i\alpha_i$ and $\PP(\text{any loss})\le\sum_i(1-\pi_i)\beta_i$ under arbitrary dependence
between steps. For i.i.d.\ probes the leading-order cost
\[
J=\textstyle\sum_i c_i\big[(1-\pi_i)\log\frac{1}{\alpha_i}/D_{0,i}+\pi_i\log\frac{1}{\beta_i}/D_{1,i}\big]
\]
is strictly convex, and subject to $\sum\alpha_i\le\alpha$, $\sum\beta_i\le\beta$ it is minimized by
\begin{equation}
\alpha_i^\star\propto w_i=\frac{c_i(1-\pi_i)}{D_{0,i}},\qquad \beta_i^\star\propto v_i=\frac{c_i\pi_i}{D_{1,i}},
\label{eq:kkt}
\end{equation}
saving $W\,\KL(\hat w\|u_k)+V\,\KL(\hat v\|u_k)$ over the uniform split ($W=\sum w_i$, $\hat w=w/W$, $u_k$ uniform; $V,\hat v$ likewise).
\end{proposition}
The guarantee part is a union bound, and it is what the saga layer is for: giving every step the full $\alpha$, the natural per-call default, multiplies the end-to-end
duplicate probability by up to $k$. The allocation part is not optimal in exact cost. It optimizes the surrogate $J$, needs priors $\pi_i$ (the per-step guarantee does not),
and is sometimes worse than uniform (Table~\ref{tab:saga}). An ambiguous compensation is asked about.

\noindent\emph{Proof.}
\emph{Guarantee.} $\PP(\dup_i)=\pi_i\PP(\dup_i\mid X_i{=}1)\le\pi_i\alpha_i$, so $\PP(\bigcup_i\dup_i)\le\sum_i\pi_i\alpha_i$ for any joint law; losses are analogous with $1-\pi_i$.
\emph{Optimum.} $J=\sum_iw_i\log\frac1{\alpha_i}+\sum_iv_i\log\frac1{\beta_i}$ is strictly convex on the open orthant and decreasing in each variable, so the constraints are active.
Stationarity, $-w_i/\alpha_i+\lambda=0$, gives $\alpha_i=w_i/\lambda$ with $\lambda=W/\alpha$; $\beta_i$ is analogous.
\emph{Saving.} $J_{\mathrm{unif}}-J^\star=\sum_iw_i\big[\log\frac k\alpha-\log\frac{W}{\alpha w_i}\big]=\sum_iw_i\log\frac{kw_i}{W}=W\,\KL(\hat w\|u_k)$, plus the analogous $\beta$ term.

\subsection{Label budgets}
\label{app:labels}
Table~\ref{tab:nsize} gives the calibration sizes that make certification possible (\S\ref{sec:design-probes}). Table~\ref{tab:cert-lagapp} extends the label budget of
Table~\ref{tab:cert-budget} to larger tied budgets and to a bootstrap lag table at the larger label count.

\begin{table}[t]
\centering
\caption{Calibration sizes that make certification possible per stratum ($\delta{=}\min(\alpha,\beta)/(10K)$, $\alpha{=}\beta$); keying by four fault classes multiplies each entry by \NsizeKeyLo{} to \NsizeKeyHi. Product: effects for two families with wrong-verdict rate $p$ to reach $\alpha/2$ (\emph{med.}: at the expected count, \emph{0}: zero wrong). Joint: effects to certify $\alpha_{\mathrm{test}}$ (\emph{0}: zero events, \emph{med.}: true rate $\alpha_{\mathrm{test}}/4$). Lag: logged times for one wait at $0.2\alpha$.}
\label{tab:nsize}
\footnotesize
\setlength{\tabcolsep}{2.2pt}
\begin{tabular}{@{}cc cc ccc c@{}}
\toprule
 & & \multicolumn{2}{c}{product} & \multicolumn{3}{c}{joint} & lag \\
\cmidrule(lr){3-4}\cmidrule(lr){5-7}
$\alpha$ & $p$ & med. & 0 & 0 & med. & 0, $k{=}3$ & 0 \\
\midrule
$10^{-2}$ & 0.5\% & 163 & 127 & 1{,}700 & 3{,}723 & 5{,}107 & 4{,}649 \\
$10^{-2}$ & 1.0\% & 193 & 127 & 1{,}700 & 3{,}723 & 5{,}107 & 4{,}649 \\
$10^{-3}$ & 0.5\% & 1{,}086 & 514 & 21{,}635 & 46{,}887 & 64{,}914 & 58{,}036 \\
$10^{-3}$ & 1.0\% & 2{,}268 & 514 & 21{,}635 & 46{,}887 & 64{,}914 & 58{,}036 \\
\bottomrule
\end{tabular}
\end{table}

% AUTO-GENERATED by artifact/scripts/render_guarantee.py -- do not edit by hand
\begin{table}[t]
\centering
\caption{\emph{Measured-parameter simulation}, Element: larger label budgets (labels / lag entries per class; ``tied'': as many lag entries as labels; ``bootstrap'': resampled from the measured lags). }
\label{tab:cert-lagapp}
\scriptsize
\setlength{\tabcolsep}{3pt}
\begin{tabular}{@{}l cc cc@{}}
\toprule
 & \multicolumn{2}{c}{product} & \multicolumn{2}{c}{joint} \\
\cmidrule(lr){2-3}\cmidrule(lr){4-5}
labels / lag entries & $P_{\mathrm v}$ (CP) & ask & $P_{\mathrm v}$ (CP) & ask \\
\midrule
10{,}000 / 10{,}000 (tied) & 0.00 (0.04) & 11.5 & 0.00 (0.04) & 2.2 \\
20{,}000 / 20{,}000 (tied) & 0.00 (0.04) & 2.2 & 0.00 (0.04) & 2.2 \\
10{,}000 / 20{,}000, bootstrap & 0.00 (0.04) & 2.5 & 0.00 (0.04) & 2.2 \\
\bottomrule
\end{tabular}
\end{table}


\subsection{Composed error bounds}
\label{app:composed}
Table~\ref{tab:cert-bound} evaluates Eq.~\eqref{eq:composed} with the budgets in place of the planner's own error rates. The $q_x$ and $\epsilon_x$ columns are rates the
simulations observed. The simulated user errs on 5\% of asks, or 10\% with a look-alike on screen, so $\epsilon_x$ depends on which effects a mode asks about. Neither
simulated world has an app that retransmits without dedup ($r_1{=}r_0{=}0$).

% AUTO-GENERATED by artifact/scripts/render_guarantee.py -- do not edit by hand
\begin{table}[t]
\centering
\caption{Composed error bounds~\eqref{eq:composed} with $r_1{=}r_0{=}0$. $q_x$: \% of effects with outcome $x$ the planner asks about; $\epsilon_x$: \% of those answers wrong. Element: \emph{measured-parameter simulation} (Table~\ref{tab:realsim}, \RsimBaseLabelN{} labels per class and outcome); default: \emph{model-based simulation}. $\alpha=\beta=1\%$; N/A: empty denominator.}
\label{tab:cert-bound}
\scriptsize
\setlength{\tabcolsep}{3pt}
\begin{tabular}{@{}l ccc ccc@{}}
\toprule
 & $q_1$ & $\epsilon_1$ & dup bound & $q_0$ & $\epsilon_0$ & loss bound \\
\midrule
Element, product & 1.92 & 10.06 & 1.19 & 4.19 & 6.01 & 1.25 \\
Element, joint & 65.06 & 6.33 & 5.12 & 100.00 & 5.63 & 6.63 \\
default, product & 22.01 & 5.99 & 2.32 & 22.99 & 7.90 & 2.82 \\
default, joint & 70.46 & 6.39 & 5.50 & 71.11 & 7.18 & 6.11 \\
\bottomrule
\end{tabular}
\end{table}


\subsection{Lag anchors}
\label{app:anchor}
Table~\ref{tab:cert-clock} compares waiting from the anchor $r_e$ of \S\ref{sec:model} with waiting from $t_0$ for every fault class, after an outage of 5 or 15\,s
before the anchor. Waiting from $r_e$ keeps every row within the budget. Waiting from $t_0$ puts every row over it, for both product and joint certification.

% AUTO-GENERATED by artifact/scripts/render_guarantee.py -- do not edit by hand
\begin{table}[t]
\centering
\caption{\emph{Measured-parameter simulation}, Element, 12{,}000 interrupted workflows per row, 5{,}000 labels per class and outcome. The app or network is down for 5 or 15\,s before the anchor $r_e$. Rows wait from the anchor $r_e$ of \S\ref{sec:model} or from $t_0$ for every class. $P_{\mathrm v}$: \% of committed effects the planner retried by itself, with the upper end of its two-sided 95\% Clopper--Pearson interval (CP). Ask: \% of workflows. Budget $\alpha=1\%$. $^\dagger$: CP bound over budget.}
\label{tab:cert-clock}
\scriptsize
\setlength{\tabcolsep}{3pt}
\begin{tabular}{@{}l cc cc@{}}
\toprule
 & \multicolumn{2}{c}{product} & \multicolumn{2}{c}{joint} \\
\cmidrule(lr){2-3}\cmidrule(lr){4-5}
wait from, downtime & $P_{\mathrm v}$ (CP) & ask & $P_{\mathrm v}$ (CP) & ask \\
\midrule
$r_e$, 5\,s & 0.00 (0.04) & 2.7 & 0.00 (0.04) & 12.7 \\
$r_e$, 15\,s & 0.00 (0.04) & 2.7 & 0.00 (0.04) & 12.7 \\
$t_0$, 5\,s & 18.82 (19.62)$^\dagger$ & 21.2 & 29.93 (30.87)$^\dagger$ & 21.6 \\
$t_0$, 15\,s & 43.77 (44.78)$^\dagger$ & 20.7 & 46.42 (47.44)$^\dagger$ & 25.8 \\
\bottomrule
\end{tabular}
\end{table}

